% F5 · Plazos concurrentes: rige el más largo (D12·05)
\documentclass[border=10pt]{standalone}
% ------------------------------------------------------------------
% Preámbulo común de las figuras del curso C1
% Paleta y tipografía alineadas con ../../estilo.tex
% ------------------------------------------------------------------
\usepackage[utf8]{inputenc}
\usepackage[T1]{fontenc}
\usepackage[spanish,es-noquoting]{babel}
\usepackage{helvet}
\renewcommand{\familydefault}{\sfdefault}
\usepackage{xcolor}
\usepackage{tikz}
\usetikzlibrary{arrows.meta,positioning,decorations.pathreplacing,calc,shapes.geometric,fit,backgrounds}

\definecolor{cHead}{RGB}{68,84,106}
\definecolor{cAcc}{RGB}{0,102,153}
\definecolor{cGris}{gray}{0.42}
\definecolor{cSuave}{RGB}{240,244,248}
\definecolor{cBien}{RGB}{34,124,68}
\definecolor{cMal}{RGB}{176,42,42}
\definecolor{cBienSuave}{RGB}{230,243,234}
\definecolor{cMalSuave}{RGB}{250,235,233}
\definecolor{cAccSuave}{RGB}{225,238,246}

\tikzset{
  caja/.style     = {rectangle, rounded corners=2.5pt, draw=cHead, line width=.7pt,
                     fill=cSuave, align=center, inner sep=6pt, font=\small},
  cajaAcc/.style  = {caja, draw=cAcc, fill=cAccSuave},
  cajaOk/.style   = {caja, draw=cBien, fill=cBienSuave},
  cajaNo/.style   = {caja, draw=cMal, fill=cMalSuave, dashed},
  cajaTenue/.style= {caja, draw=cGris, fill=white},
  rotulo/.style   = {font=\scriptsize\itshape, text=cGris, align=center},
  etiq/.style     = {font=\scriptsize, text=cHead, align=center,
                     fill=white, inner sep=1.5pt},
  fl/.style       = {-{Stealth[length=4.5pt]}, draw=cHead, line width=.7pt},
  flAcc/.style    = {fl, draw=cAcc},
  flNo/.style     = {-{Stealth[length=4.5pt]}, draw=cMal, line width=.7pt, dashed},
  llave/.style    = {decorate, decoration={brace, amplitude=4pt, raise=2pt},
                     draw=cGris, line width=.6pt},
  titulo/.style   = {font=\small\bfseries, text=cHead},
}

% Cabecera de figura: #1 número, #2 título
\newcommand{\figtit}[2]{\node[font=\footnotesize\bfseries, text=cGris, anchor=west]}

\begin{document}
\begin{tikzpicture}[x=0.48cm, y=1cm]

% ---------------- barras, con el rótulo encima ----------------------
% #1 y  #2 hasta  #3 color trazo  #4 relleno  #5 rótulo
\newcommand{\barra}[5]{%
  \node[font=\scriptsize, text=cHead, anchor=west] at (0,#1+.34) {#5};
  \fill[#4] (0,#1-.17) rectangle (#2,#1+.17);
  \draw[#3, line width=.8pt] (0,#1-.17) rectangle (#2,#1+.17);}

\barra{4.15}{25}{cMal}{cMalSuave}
  {\textbf{25 años} · archivo maestro del ensayo clínico, desde su finalización}
\barra{2.95}{15}{cAcc}{cAccSuave}
  {\textbf{15 años} · conservación comprometida en el plan}
\barra{1.75}{3}{cHead}{cSuave}
  {\textbf{3 años} · límite de la autorización de acceso a datos preexistentes}

% historia clínica: obligación ajena y con otro origen
\node[font=\scriptsize, text=cGris, anchor=west] at (0,0.55+.34)
  {\textbf{mín. 5 años} · historia clínica \emph{(del centro, desde el alta de cada proceso: otro origen)}};
\draw[cGris, line width=.8pt, dash pattern=on 2.5pt off 2pt]
  (0,0.55-.17) rectangle (5,0.55+.17);

% ---------------- eje ------------------------------------------------
\draw[cGris, line width=.6pt, -{Stealth[length=5pt]}] (0,-.35) -- (28,-.35);
\foreach \x/\et in {0/0, 5/5, 10/10, 15/15, 20/20, 25/25}
  \draw[cGris] (\x,-.27) -- (\x,-.43) node[below, font=\scriptsize, text=cGris] {\et};
\node[font=\scriptsize\itshape, text=cGris, anchor=west] at (0,-.95)
  {años desde el cierre del proyecto};

% ---------------- lo que rige ----------------------------------------
\draw[cMal, line width=1pt, dash pattern=on 3pt off 2pt] (25,4.75) -- (25,-.35);
\node[font=\small\bfseries, text=cMal, anchor=south] at (25,4.82) {rige el más largo};

% ---------------- y después ------------------------------------------
\node[cajaTenue, font=\scriptsize, anchor=west, text width=3.3cm] (bloq) at (26.2,3.4)
  {\textbf{Bloqueo}\\ mientras dure la prescripción};
\node[cajaTenue, font=\scriptsize, anchor=west, text width=3.3cm] (borr) at (26.2,1.6)
  {\textbf{Borrado seguro}\\ con acta de destrucción};
\draw[fl] (25,3.4) -- (bloq.west);
\draw[fl] (bloq.south) -- (borr.north);

\node[rotulo, text width=19cm, align=left, anchor=north west] at (0,-1.35)
  {Cumplir el plazo corto no exime del largo. La obligación sobre la historia clínica es del centro,
   no del proyecto, y se cuenta desde otro momento: lo que el proyecto conserva es su copia extraída,
   con su propio régimen. Los metadatos, la documentación y los datos agregados no se destruyen.};

\end{tikzpicture}
\end{document}
